% Bagean: metodhe
% Bab: induksi
% Pérangan: induksi-struktural

\documentclass[../../../include/open-logic-section]{subfiles}

\begin{document}

\olfileid{mth}{ind}{sti}

\olsection{Induksi Struktural}

Nganti saiki induksi digunakaké kanggo
mbuktekaké pratelan tumrap kabèh bilangan
asli. Nanging prinsip kang padha bisa
langsung digunakaké kanggo mbuktekaké
pratelan tumrap kabèh !!{element}s
saka himpunan kang didhéfinisikaké
kanthi induktif. Iki kerep diarani
induksi \emph{struktural}, amarga
gumantung marang struktur obyèk
kang didhéfinisikaké kanthi induktif.

Lumrahé dhéfinisi induktif mènèhi
(a)~dhaptar !!{element}s ``wiwitan''
saka himpunan, lan (b)~dhaptar
operasi kang nggawe !!{element}s
anyar saka kang wis ana. Ing
conto term trep, obyèk wiwitan
yaiku aksara. Operasiné mung siji:
\begin{align*}
  o(s_1, s_2) = & [s_1 \circ s_2]
\end{align*}
Bilangan asli~$\Nat$ dhéwé uga
bisa dianggep didhéfinisikaké
kanthi induktif: obyèk wiwitan
yaiku~$0$, lan operasiné
fungsi panerus~$x + 1$.

Kanggo mbuktekaké yèn kabèh
anggota himpunan induktif,
yaiku saben !!{element} ing
himpunan mau, nduwèni
sipat~$P$, kita kudu:
\begin{enumerate}
\item Mbuktekaké yèn obyèk
  wiwitan nduwèni~$P$.
\item Mbuktekaké yèn kanggo
  saben operasi~$o$, yèn
  argumèné nduwèni~$P$,
  asilé uga nduwèni sipat iku.
\end{enumerate}
Upamané, kanggo mbuktekaké
pratelan tumrap kabèh term
trep, kita buktèkaké kanggo
kabèh aksara, lan kanggo
$[s_1 \circ s_2]$ kanthi
anggepan yèn pratelan iku
bener tumrap $s_1$ lan
$s_2$ siji-siji.

\begin{prop}
  Cacah $[$ padha karo cacah
  $]$ ing sembarang term
  trep~$t$.
\end{prop}

\begin{proof}
Kita migunakaké induksi
struktural. Term trep
didhéfinisikaké kanthi
induktif: aksara minangka
obyèk wiwitan lan operasi
$o$ kanggo mbangun term
trep anyar saka kang wis ana.
\begin{enumerate}
\item Pratelan bener kanggo
  saben aksara, amarga
  cacah $[$ ing aksara
  tunggal yaiku~$0$ lan
  cacah $]$ uga~$0$.
\item Anggepen cacah $[$
  ing $s_1$ padha karo
  cacah $]$, lan bab kang
  padha bener kanggo $s_2$.
  Cacah $[$ ing
  $o(s_1, s_2)$, yaiku
  ing $[s_1 \circ s_2]$,
  padha karo jumlah cacah
  $[$ ing $s_1$ lan
  $s_2$ ditambah siji.
  Cacah $]$ ing
  $o(s_1, s_2)$ padha
  karo jumlah cacah $]$
  ing $s_1$ lan $s_2$
  ditambah siji. Mula
  cacah $[$ ing
  $o(s_1, s_2)$ padha
  karo cacah $]$ ing
  $o(s_1,s_2)$.
\end{enumerate}
\end{proof}

\begin{prob}
  Buktèkna kanthi induksi
  struktural yèn ora ana
  term trep kang diwiwiti
  déning~$]$.
\end{prob}

Ayo deleng bukti induksi
struktural liyané: bagéan
wiwitan sejati saka rerangken
simbol~$t$ yaiku rerangken
ora kosong~$s$ kang saben
simbolé padha karo simbol
wiwitan~$t$ yèn diwaca
saka kiwa, nanging $t$
luwih dawa. Upamané,
$[a \circ {}$ iku bagéan
wiwitan sejati saka
$[a \circ b]$, nanging
$[b \circ {}$ dudu
(béda ing simbol kapindho),
lan $[a
\circ b]$ uga dudu
(dawané padha).

\begin{prop}\ollabel{prop:initial}
  Saben bagéan wiwitan sejati
  kang ora kosong saka term
  trep~$t$ ngemot $[$ luwih
  akèh tinimbang $]$.
\end{prop}

\begin{proof}
  Kanthi induksi ing~$t$:
  \begin{enumerate}
  \item Yèn $t$ mung aksara,
    $t$ ora nduwèni bagéan
    wiwitan sejati kang ora kosong.
  \item Yèn $t = [s_1 \circ s_2]$
    kanggo term trep $s_1$
    lan~$s_2$, lan $r$ bagéan
    wiwitan sejati kang ora kosong
    saka $t$, ana sawatara
    kamungkinan:
    \begin{enumerate}
    \item Yèn $r$ mung $[$,
      $r$ nduwèni cacah $[$
      luwih siji tinimbang~$]$.
    \item Yèn $r$ awujud
      $[r_1$ kanthi $r_1$
      bagéan wiwitan sejati
      saka~$s_1$, amarga
      $s_1$ term trep,
      miturut hipotesis
      induksi $r_1$
      ngemot $[$ luwih
      akèh tinimbang $]$,
      lan $[r_1$ uga
      mangkono.
    \item Yèn $r$ awujud
      $[s_1$ utawa
      $[s_1 \circ {}$,
      miturut asil sadurungé
      cacah $[$ lan $]$
      ing~$s_1$ padha.
      Dadi cacah $[$
      ing $[s_1$ utawa
      $[s_1 \circ {}$
      luwih siji tinimbang
      cacah~$]$.
    \item Yèn $r$ awujud
      $[s_1 \circ r_2$
      kanthi $r_2$ bagéan
      wiwitan sejati
      saka~$s_2$, miturut
      hipotesis induksi
      $r_2$ ngemot $[$
      luwih akèh tinimbang
      $]$. Miturut asil
      sadurungé, cacah $[$
      lan $]$ ing~$s_1$
      padha. Mula cacah $[$
      ing $[s_1 \circ r_2$
      luwih akèh tinimbang
      cacah~$]$.
    \item Yèn $r$ awujud
      $[s_1 \circ s_2$,
      miturut asil sadurungé
      cacah $[$ lan $]$
      ing $s_1$ padha,
      lan semono uga
      ing~$s_2$. Dadi
      cacah $[$ ing
      $[s_1 \circ s_2$
      luwih siji tinimbang
      cacah~$]$.
    \end{enumerate}
  \end{enumerate}
\end{proof}

\end{document}
